% =====================================================================
%  Chương 4 — Mạng nơ-ron hồi quy (RNN)
%  Số liệu: \SL{ch4_*} từ src/TieuLuan 03/outputs/metrics/ch4_*.json; mã: \maNguon từ tl3/ch4_*, scripts/ch4_*.
% =====================================================================
\chapter{Mạng nơ-ron hồi quy (Recurrent Neural Networks -- RNN)}
\label{chap:rnn}

Chương~3 cho thấy mạng tích chập khai thác cấu trúc \emph{không gian} của ảnh. Chương này chuyển sang dữ liệu có
cấu trúc \emph{thời gian}: doanh thu từng ngày, giá cổ phiếu từng phiên, giá vàng, câu văn, tín hiệu tiếng nói. Ở các
dữ liệu đó thứ tự là một phần của thông tin, và giá trị hôm nay chỉ có nghĩa khi đặt cạnh những gì đã xảy ra trước nó.
Phần lý thuyết (mục~\ref{sec:rnn-khainiem}–\ref{sec:rnn-sieuthamso}) đi từ khái niệm, lịch sử, công thức của RNN đơn
giản đến lan truyền ngược qua thời gian, hiện tượng gradient biến mất, các kiến trúc có cổng LSTM và GRU, rồi toàn bộ
quy trình dự báo chuỗi thời gian. Mục~\ref{sec:rnn-scratch} cài đặt RNN từ đầu bằng NumPy và kiểm tra gradient bằng sai
phân số. Ba thực nghiệm (mục~\ref{sec:rnn-tn1}–\ref{sec:rnn-tn3}) dự báo doanh thu thương mại điện tử, giá cổ phiếu
AAPL và giá quỹ vàng GLD bằng SimpleRNN, LSTM, GRU và BiLSTM trên cả PyTorch lẫn Keras; mỗi thực nghiệm kết thúc bằng
phần triển khai thành ứng dụng web.

% =====================================================================
\section{Khái niệm về RNN}
\label{sec:rnn-khainiem}

\begin{dinhnghia}{Định nghĩa}
Mạng nơ-ron hồi quy (recurrent neural network, RNN) là mạng nơ-ron xử lý một chuỗi $x_1, x_2, \dots, x_T$ \emph{lần
lượt từng phần tử}, đồng thời duy trì một \emph{trạng thái ẩn} (hidden state) $h_t$ tóm tắt những gì mạng đã đọc tới
bước $t$. Ở mỗi bước, trạng thái mới được tính từ đầu vào hiện tại và trạng thái cũ bằng \emph{cùng một} hàm có tham
số dùng chung:
\[ h_t = f_\theta(h_{t-1}, x_t), \qquad \hat y_t = g_\theta(h_t). \]
\end{dinhnghia}

Khác với mạng truyền thẳng (feed-forward) coi mỗi mẫu là một vector cố định, RNN có một vòng phản hồi: đầu ra của
nơ-ron ẩn ở bước trước được đưa trở lại làm đầu vào ở bước sau. Nhờ vòng phản hồi, RNN có ba tính chất mà mạng truyền
thẳng không có:
\begin{itemize}
  \item \textbf{Nhận chuỗi độ dài bất kỳ.} Số tham số không phụ thuộc $T$ vì cùng một bộ trọng số được áp dụng ở mọi bước.
  \item \textbf{Có bộ nhớ.} $h_t$ là hàm của toàn bộ $x_1,\dots,x_t$, nên quyết định ở bước $t$ có thể dựa vào thông tin
        xuất hiện rất lâu trước đó (ít nhất về nguyên tắc).
  \item \textbf{Bất biến theo dịch chuyển thời gian.} Một mẫu hình (ví dụ ``giá tăng ba phiên liên tiếp'') được nhận
        ra bằng cùng một phép tính dù nó xảy ra ở đầu hay cuối chuỗi, tương tự như tích chập dùng chung bộ lọc trên
        mọi vị trí của ảnh.
\end{itemize}
Có thể xem RNN là một \emph{hệ động lực rời rạc} có đầu vào: $h_t$ là trạng thái, $x_t$ là tín hiệu điều khiển, và
mạng học chính hàm chuyển trạng thái $f_\theta$ từ dữ liệu \cite{ch4:goodfellow2016}.

% =====================================================================
\section{Lịch sử phát triển của RNN}
\label{sec:rnn-lichsu}

\hinhtikz{ch4_history}{Các mốc chính trong lịch sử mạng hồi quy, từ mạng Hopfield (1982) đến các mô hình không gian
trạng thái chọn lọc (2023).}{ch4-history}

Lịch sử RNN gắn với câu hỏi làm sao để mạng nơ-ron có trí nhớ (Hình~\ref{fig:ch4-history}):
\begin{itemize}
  \item \textbf{1982 -- Mạng Hopfield} \cite{ch4:hopfield1982}: mạng hồi tiếp toàn phần dùng như bộ nhớ liên tưởng,
        hội tụ về các điểm cực tiểu năng lượng. Đây chưa phải mô hình chuỗi nhưng đặt nền cho ý tưởng trạng thái lưu
        thông tin qua phản hồi.
  \item \textbf{1986 -- Mạng Jordan và lan truyền ngược}: Jordan \cite{ch4:jordan1986} đưa đầu ra trở lại làm ngữ cảnh;
        cùng năm, Rumelhart, Hinton và Williams phổ biến lan truyền ngược \cite{ch4:rumelhart1986}, kèm cách ``mở''
        mạng hồi tiếp theo thời gian để huấn luyện.
  \item \textbf{1990 -- Mạng Elman (Simple Recurrent Network)} \cite{ch4:elman1990}: trạng thái ẩn được sao chép thành
        ``đơn vị ngữ cảnh'' cho bước sau; đây chính là SimpleRNN dùng trong chương này. Cùng năm, Werbos hệ thống hoá
        BPTT \cite{ch4:werbos1990}.
  \item \textbf{1994 -- Phát hiện gradient biến mất}: Bengio, Simard và Frasconi \cite{ch4:bengio1994} chứng minh huấn
        luyện bằng gradient không học được phụ thuộc dài hạn một cách ổn định.
  \item \textbf{1997 -- LSTM và RNN hai chiều}: Hochreiter và Schmidhuber đề xuất LSTM với ``băng chuyền lỗi không
        đổi'' \cite{ch4:hochreiter1997}; Schuster và Paliwal đề xuất BiRNN \cite{ch4:schuster1997}. Năm 2000, Gers và
        cộng sự thêm cổng quên \cite{ch4:gers2000}, cho ra LSTM như ngày nay.
  \item \textbf{2013--2015 -- Thời kỳ hoàng kim}: RNN sâu đạt kỷ lục nhận dạng tiếng nói \cite{ch4:graves2013}; GRU
        \cite{ch4:cho2014} và kiến trúc seq2seq \cite{ch4:sutskever2014} đưa dịch máy nơ-ron vào thực tế; cơ chế chú ý
        (attention) \cite{ch4:bahdanau2015} gỡ nút thắt vector ngữ cảnh cố định.
  \item \textbf{2017 -- Transformer} \cite{ch4:vaswani2017}: bỏ hẳn hồi tiếp, chỉ dùng attention nên huấn luyện song song
        được; Transformer thay RNN ở hầu hết bài toán ngôn ngữ quy mô lớn.
  \item \textbf{2023 -- Hồi tiếp trở lại}: các mô hình không gian trạng thái chọn lọc như Mamba \cite{ch4:gu2023} và
        RWKV kết hợp suy luận tuyến tính theo độ dài chuỗi của RNN với khả năng huấn luyện song song, cho thấy ý tưởng
        trạng thái hồi tiếp vẫn còn giá trị.
\end{itemize}

% =====================================================================
\section{Cấu trúc và nguyên lý hoạt động của RNN}
\label{sec:rnn-cautruc}

\subsection{Cấu trúc của RNN}

RNN Elman gồm ba nhóm trọng số: $W_{xh}\in\mathbb{R}^{H\times F}$ nối đầu vào với lớp ẩn, $W_{hh}\in\mathbb{R}^{H\times H}$
nối lớp ẩn với chính nó ở bước trước, và $W_{hy}\in\mathbb{R}^{O\times H}$ nối lớp ẩn với đầu ra. Với $F$ đặc trưng
đầu vào, $H$ đơn vị ẩn và $O$ đầu ra:
\begin{align}
  a_t &= W_{xh}\,x_t + W_{hh}\,h_{t-1} + b_h, \label{eq:rnn-a}\\
  h_t &= \tanh(a_t), \label{eq:rnn-h}\\
  \hat y_t &= W_{hy}\,h_t + b_y. \label{eq:rnn-y}
\end{align}
Số tham số của lớp hồi tiếp là $H(F+H+1)$, không phụ thuộc độ dài chuỗi. Hình~\ref{fig:ch4-unfold} vẽ cùng một mạng
ở hai dạng: dạng gập với vòng phản hồi $W_{hh}$, và dạng \emph{mở theo thời gian} (unfold) trong đó mỗi bước là một
bản sao của cùng một ô. Dạng mở biến RNN thành một mạng truyền thẳng rất sâu (sâu $T$ tầng) có trọng số dùng chung,
và đó là cách huấn luyện nó.

\hinhtikz{ch4_unfold}{RNN ở dạng gập (trái) và dạng mở theo thời gian (phải): mọi bước dùng chung $W_{xh}$, $W_{hh}$,
$W_{hy}$.}{ch4-unfold}

\hinhtikz{ch4_cell}{Bên trong một ô SimpleRNN: hai phép nhân ma trận, cộng bias, hàm $\tanh$.}{ch4-cell}

\subsection{Nguyên lý hoạt động}
\label{sec:rnn-nguyenly}

Hình~\ref{fig:ch4-cell} cho thấy ở mỗi bước, ô RNN thực hiện bốn thao tác: (1) chiếu đầu vào $x_t$ vào không gian ẩn
bằng $W_{xh}$; (2) chiếu trạng thái cũ $h_{t-1}$ bằng $W_{hh}$; (3) cộng hai kết quả với bias thành tổng kích hoạt
$a_t$; (4) nén $a_t$ vào khoảng $(-1,1)$ bằng $\tanh$ để được $h_t$. Trạng thái $h_t$ vừa là ``bộ nhớ'' chuyển cho bước
sau, vừa là biểu diễn để lớp đầu ra dự báo. Trạng thái ban đầu thường đặt $h_0=\mathbf 0$. Với bài toán dự báo
many-to-one của chương này, chỉ trạng thái cuối $h_T$ được đưa qua lớp tuyến tính để sinh một giá trị.

\subsection{Ưu điểm}
\begin{itemize}
  \item Xử lý chuỗi độ dài thay đổi với số tham số cố định; mô hình gọn (SimpleRNN 64 đơn vị cho AAPL chỉ có
        \SL{ch4_aapl_torch/rnn/params} tham số).
  \item Chia sẻ tham số theo thời gian giúp tổng quát hoá: mẫu hình học ở đoạn này dùng được ở đoạn khác.
  \item Suy luận trực tuyến (streaming): mỗi bước chỉ cần $h_{t-1}$ và $x_t$, chi phí $O(1)$ theo độ dài lịch sử, phù
        hợp cảm biến, tài chính tần suất cao.
  \item Mô hình hoá tự nhiên tính nhân quả: $h_t$ chỉ phụ thuộc quá khứ.
\end{itemize}

\subsection{Nhược điểm}
\begin{itemize}
  \item \textbf{Phụ thuộc dài hạn khó học} do gradient biến mất hoặc bùng nổ (mục~\ref{sec:rnn-vanish}).
  \item \textbf{Tuần tự, khó song song}: bước $t$ phải đợi bước $t-1$, nên huấn luyện trên GPU kém hiệu quả hơn CNN
        hay Transformer với chuỗi dài.
  \item \textbf{Nút thắt thông tin}: toàn bộ lịch sử bị nén vào vector $h_t$ kích thước cố định.
  \item Nhạy với khởi tạo, learning rate và chuẩn hoá dữ liệu; dễ quá khớp trên chuỗi ngắn.
\end{itemize}

\subsection{Ứng dụng}
RNN và các biến thể từng là kiến trúc chủ đạo cho: nhận dạng tiếng nói và chữ viết tay \cite{ch4:graves2013}; dịch
máy và tóm tắt văn bản (seq2seq) \cite{ch4:sutskever2014}; mô hình ngôn ngữ, gõ dự đoán trên bàn phím điện thoại;
phân tích cảm xúc; dự báo chuỗi thời gian (nhu cầu điện, doanh số, giá tài sản, lưu lượng giao thông); phát hiện bất
thường trong log và cảm biến IoT; sinh nhạc, sinh chuỗi hoạt động của người dùng. Trong chương này RNN được dùng cho
bài toán dự báo chuỗi thời gian đa biến.

% =====================================================================
\section{Biểu diễn dữ liệu chuỗi thời gian}
\label{sec:rnn-bieudien}

Một chuỗi thời gian đa biến là ma trận $\mathbf{S}\in\mathbb{R}^{N\times F}$: $N$ bước thời gian (ngày, phiên),
mỗi bước $F$ đặc trưng. Để huấn luyện theo lô, ta cắt $\mathbf S$ thành các cửa sổ độ dài $T$ và xếp thành tensor ba
chiều (Hình~\ref{fig:ch4-tensor}). Bảng~\ref{tab:ch4-repr} tóm tắt quy ước dùng trong chương.

\hinhtikz{ch4_tensor}{Tensor đầu vào của RNN: lô $B$ cửa sổ, mỗi cửa sổ $T$ bước, mỗi bước $F$ đặc trưng.}{ch4-tensor}

\begin{table}[H]\centering\caption{Quy ước biểu diễn chuỗi trong chương}\label{tab:ch4-repr}
\small
\begin{tabular}{L{3.6cm}L{4.2cm}L{6.2cm}}
\toprule
Ký hiệu & Kích thước & Ý nghĩa \\
\midrule
$x_t$ & $(F,)$ & Vector đặc trưng tại bước $t$ (ví dụ open, high, low, close, adjclose, volume) \\
$X_i$ & $(T, F)$ & Một cửa sổ $T=30$ bước liên tiếp \\
$\mathbf X$ & $(B, T, F)$ & Một lô cửa sổ đưa vào mạng; e-commerce dùng $(\cdot,30,9)$, AAPL $(\cdot,30,6)$, vàng $(\cdot,30,13)$ \\
$y_i$ & vô hướng & Giá trị cần dự báo ở bước ngay sau cửa sổ: Total\_Amount hoặc close \\
$h_t$ & $(H,)$ & Trạng thái ẩn, $H=64$ trong mọi thực nghiệm \\
\bottomrule
\end{tabular}
\end{table}

% =====================================================================
\section{Hàm kích hoạt trong RNN}
\label{sec:rnn-kichhoat}

Hàm $\tanh$ là lựa chọn mặc định cho trạng thái ẩn vì đầu ra nằm trong $(-1,1)$, đối xứng quanh 0 và bão hoà nên
chặn được trạng thái khỏi phình to qua nhiều bước. Sigmoid có đầu ra $(0,1)$ nên phù hợp làm ``van'' (cổng) trong
LSTM/GRU: giá trị 0 là đóng, 1 là mở. ReLU ít được dùng làm kích hoạt trạng thái của RNN truyền thống vì không bị chặn
trên; mục~\ref{sec:rnn-vanish} cho thấy RNN-ReLU với bán kính phổ $\rho(W_{hh})>1$ làm gradient bùng nổ hàng chục triệu
lần.

\begin{table}[H]\centering\caption{Các hàm kích hoạt thường gặp trong RNN}\label{tab:ch4-act}
\small
\begin{tabular}{L{2.2cm}L{4.3cm}L{3.2cm}L{4.5cm}}
\toprule
Hàm & Công thức & Đạo hàm & Vai trò trong RNN \\
\midrule
Sigmoid & $\sigma(x)=\dfrac{1}{1+e^{-x}}$ & $\sigma(1-\sigma)\le\tfrac14$ & Cổng quên/vào/ra của LSTM, GRU \\[6pt]
tanh & $\tanh x=\dfrac{e^{x}-e^{-x}}{e^{x}+e^{-x}}$ & $1-\tanh^2 x\le 1$ & Trạng thái ẩn, ứng viên cell \\[6pt]
ReLU & $\max(0,x)$ & $\mathbf{1}[x>0]$ & Ít dùng cho trạng thái (dễ bùng nổ) \\[6pt]
Softmax & $\dfrac{e^{x_i}}{\sum_j e^{x_j}}$ & $p_i(\delta_{ij}-p_j)$ & Đầu ra phân loại, trọng số attention \\
\bottomrule
\end{tabular}
\end{table}

% =====================================================================
\section{RNN xử lý chuỗi như thế nào?}
\label{sec:rnn-xuly}

Với một cửa sổ $X=[x_1,\dots,x_T]$, lan truyền xuôi gồm các bước:
\begin{enumerate}
  \item Khởi tạo $h_0=\mathbf 0$.
  \item Với $t=1$: tính $h_1=\tanh(W_{xh}x_1+W_{hh}h_0+b_h)$.
  \item Chuyển $h_1$ sang bước 2, kết hợp với $x_2$ để tạo $h_2$; lặp tới $t=T$.
  \item Many-to-one: dùng $h_T$ để dự báo $\hat y=W_{hy}h_T+b_y$. Many-to-many: dùng mọi $h_t$.
  \item Tính mất mát, lan truyền ngược qua toàn bộ $T$ bước (BPTT) và cập nhật trọng số.
\end{enumerate}
Đoạn mã~\ref{code:ch4_scratch__RNNScratch.forward} ở mục~\ref{sec:rnn-scratch} cài đặt đúng năm bước này bằng NumPy
trong chưa đầy mười dòng.

% =====================================================================
\section{Các kiểu bài toán RNN theo quan hệ input--output}
\label{sec:rnn-iotypes}

\hinhtikz{ch4_io_types}{Năm kiểu quan hệ đầu vào -- đầu ra của mạng hồi quy; ba thực nghiệm trong chương thuộc kiểu
many-to-one (khung đỏ).}{ch4-io}

\begin{table}[H]\centering\caption{Các kiểu bài toán RNN}\label{tab:ch4-io}
\small
\begin{tabular}{L{3.4cm}L{5.2cm}L{5.6cm}}
\toprule
Kiểu & Mô tả & Ví dụ \\
\midrule
One-to-one & Một vào, một ra; thực chất là mạng thường & Phân loại ảnh đơn \\
One-to-many & Một vào, sinh chuỗi ra & Sinh chú thích ảnh, sinh nhạc từ một nốt \\
Many-to-one & Đọc cả chuỗi, ra một giá trị & Dự báo doanh thu/giá ngày kế tiếp, phân loại cảm xúc \\
Many-to-many đồng bộ & Mỗi bước vào có một bước ra & Gán nhãn từ loại, phát hiện bất thường từng thời điểm \\
Many-to-many bất đồng bộ & Encoder đọc hết rồi decoder sinh chuỗi độ dài khác & Dịch máy, dự báo nhiều bước trực tiếp \\
\bottomrule
\end{tabular}
\end{table}

% =====================================================================
\section{BPTT -- Backpropagation Through Time}
\label{sec:rnn-bptt}

Huấn luyện RNN bằng cách mở mạng theo thời gian (Hình~\ref{fig:ch4-unfold}) rồi áp dụng lan truyền ngược trên đồ thị
tính toán đã mở. Với tổng mất mát $L=\sum_t L_t$ và ký hiệu $\delta_t=\partial L/\partial h_t$, quy tắc chuỗi cho:
\begin{align}
  \delta_T &= W_{hy}^\top\,\frac{\partial L_T}{\partial \hat y_T}, \\
  \delta_{t-1} &= \underbrace{W_{hh}^\top\,\mathrm{diag}\!\big(1-h_t^2\big)}_{(\partial h_t/\partial h_{t-1})^\top}\,\delta_t
                 \;+\; W_{hy}^\top\frac{\partial L_{t-1}}{\partial\hat y_{t-1}}, \label{eq:bptt-delta}\\
  \frac{\partial L}{\partial W_{hh}} &= \sum_{t=1}^{T}\Big[(1-h_t^2)\odot\delta_t\Big]\,h_{t-1}^\top, \qquad
  \frac{\partial L}{\partial W_{xh}} = \sum_{t=1}^{T}\Big[(1-h_t^2)\odot\delta_t\Big]\,x_t^\top. \label{eq:bptt-w}
\end{align}
Hai điểm then chốt: (i) vì trọng số dùng chung, gradient của $W_{hh}$ là \emph{tổng} đóng góp từ mọi bước; (ii) tín
hiệu lỗi từ bước $T$ về bước $t$ phải đi qua tích của $T-t$ ma trận Jacobian
\begin{equation}
  \frac{\partial h_T}{\partial h_t}=\prod_{k=t+1}^{T}\frac{\partial h_k}{\partial h_{k-1}}
  =\prod_{k=t+1}^{T}\mathrm{diag}\big(1-h_k^2\big)\,W_{hh}. \label{eq:jacobian}
\end{equation}
Trong thực tế, chuỗi rất dài được cắt thành đoạn (truncated BPTT); ở chương này mỗi cửa sổ chỉ 30 bước nên BPTT đầy
đủ được dùng.

\hinhtikz{ch4_bptt}{Lan truyền ngược qua thời gian: lỗi từ mỗi $L_t$ đi ngược về các bước trước qua chuỗi Jacobian
$\partial h_t/\partial h_{t-1}$ (mũi tên đỏ).}{ch4-bptt}

% =====================================================================
\section{Vanishing gradient và exploding gradient}
\label{sec:rnn-vanish}

Lấy chuẩn hai vế của~\eqref{eq:jacobian}: vì $|1-h_k^2|\le1$,
\begin{equation}
  \Big\|\frac{\partial h_T}{\partial h_t}\Big\| \le \prod_{k=t+1}^{T}\big\|\mathrm{diag}(1-h_k^2)\big\|\,\|W_{hh}\|
  \le \|W_{hh}\|^{\,T-t}.
\end{equation}
Nếu giá trị kỳ dị lớn nhất của $W_{hh}$ nhỏ hơn 1 (hoặc $\tanh$ bão hoà), tích co theo cấp số nhân và
\textbf{gradient biến mất}: bước xa không còn ảnh hưởng tới cập nhật, mạng ``quên'' phụ thuộc dài hạn. Ngược lại, nếu
bán kính phổ $\rho(W_{hh})>1$ và kích hoạt không bão hoà, tích phình to và \textbf{gradient bùng nổ}: loss nhảy vọt,
trọng số thành NaN \cite{ch4:bengio1994,ch4:pascanu2013}.

\hinhtikz{ch4_vanish}{Trực giác về gradient biến mất (trái) và bùng nổ (phải): độ lớn gradient theo số bước lan truyền
ngược, cùng các giải pháp tương ứng.}{ch4-vanish}

\paragraph{Đo trực tiếp.} Để thấy hiện tượng bằng số, thí nghiệm \texttt{gradient\_flow}
(Đoạn mã~\ref{code:ch4_extras__gradient_flow}) cho chuỗi ngẫu nhiên 100 bước qua RNN, LSTM, GRU khởi tạo ngẫu nhiên
(ẩn 64, trung bình 20 lần khởi tạo) và đo $\|\partial h_T/\partial x_t\|$ theo khoảng cách $T-t$.

\maDai{ch4_extras__gradient_flow}{Đo chuẩn gradient $\|\partial h_T/\partial x_t\|$ theo khoảng cách thời gian bằng
autograd của PyTorch}

\hinhpdf{ch4_gradflow}{1}{Chuẩn gradient theo khoảng cách thời gian (thang log). Trái: khởi tạo mặc định của PyTorch và
LSTM có bias cổng quên bằng 3. Phải: RNN-ReLU khi đặt bán kính phổ của $W_{hh}$.}{ch4-gradflow}
\bang{ch4_gradflow}{Chuẩn $\|\partial h_T/\partial x_t\|$ tại một số khoảng cách $T-t$}{ch4-gradflow}

\nhanxet{Với SimpleRNN, gradient từ bước cách 50 bước chỉ còn \SL{ch4_extras/gradient/summary/rnn/k50} và tới 99
bước còn \SL{ch4_extras/gradient/summary/rnn/k99}: thực tế bằng 0. Đáng chú ý, LSTM và GRU \emph{lúc khởi tạo} cũng suy
giảm gần như vậy (\SL{ch4_extras/gradient/summary/lstm/k50} ở 50 bước), vì bias cổng quên bằng 0 làm
$f_t=\sigma(\cdot)\approx0{,}5$ và cell state mất một nửa mỗi bước. Chỉ cần đặt bias cổng quên bằng 3
($f_t\approx0{,}95$), chuẩn gradient sau 50 bước vẫn là \SL{ch4_extras/gradient/summary/lstm_fb3/k50} và sau 99 bước
là \SL{ch4_extras/gradient/summary/lstm_fb3/k99}: lớn hơn SimpleRNN khoảng $10^{20}$ lần. Lợi thế của LSTM vì vậy
không tự có, mà đến từ việc cổng quên \emph{học được} (hoặc được khởi tạo) giá trị gần 1 \cite{ch4:gers2000}. Bên
phải, RNN-ReLU với $\rho=1{,}5$ cho gradient tăng tới \SL{ch4_extras/gradient/summary/lin_rho1.5/k99} sau 99 bước:
đúng dạng bùng nổ cấp số nhân $\rho^{T-t}$.}

\begin{table}[H]\centering\caption{Hai vấn đề gradient của RNN và giải pháp}\label{tab:ch4-vanish}
\small
\begin{tabular}{L{2.4cm}L{3.8cm}L{3.3cm}L{4.7cm}}
\toprule
Vấn đề & Cơ chế & Hậu quả & Giải pháp thường gặp \\
\midrule
Gradient biến mất & Tích nhiều Jacobian có chuẩn $<1$ & Không học được phụ thuộc dài hạn & LSTM/GRU, bias cổng quên dương, khởi tạo trực giao, skip connection \\
Gradient bùng nổ & Tích nhiều Jacobian có chuẩn $>1$ & Loss dao động, NaN & Cắt chuẩn gradient (clipping), learning rate nhỏ, kích hoạt bão hoà \\
\bottomrule
\end{tabular}
\end{table}

Cắt chuẩn gradient (gradient clipping) \cite{ch4:pascanu2013} giới hạn độ dài vector gradient mà không đổi hướng:
$g\leftarrow g\cdot\min\!\big(1, c/\|g\|\big)$. Mọi mô hình của chương dùng $c=1$
(\texttt{clip\_grad\_norm\_} trong PyTorch, \texttt{clipnorm=1.0} trong Keras).

% =====================================================================
\section{LSTM -- Long Short-Term Memory}
\label{sec:rnn-lstm}

LSTM \cite{ch4:hochreiter1997,ch4:gers2000} thêm một \emph{cell state} $c_t$ làm bộ nhớ dài hạn, chạy dọc theo chuỗi
như một băng chuyền và chỉ bị sửa đổi bằng phép nhân với cổng và phép cộng (Hình~\ref{fig:ch4-lstm}). Ba cổng sigmoid
điều khiển dòng thông tin:
\begin{align}
  f_t &= \sigma\big(W_f[h_{t-1},x_t]+b_f\big) && \text{cổng quên: giữ bao nhiêu của } c_{t-1}\\
  i_t &= \sigma\big(W_i[h_{t-1},x_t]+b_i\big) && \text{cổng vào: ghi bao nhiêu thông tin mới}\\
  \tilde c_t &= \tanh\big(W_c[h_{t-1},x_t]+b_c\big) && \text{nội dung ứng viên}\\
  c_t &= f_t\odot c_{t-1} + i_t\odot\tilde c_t && \text{cập nhật cộng tính}\\
  o_t &= \sigma\big(W_o[h_{t-1},x_t]+b_o\big) && \text{cổng ra}\\
  h_t &= o_t\odot\tanh(c_t).
\end{align}
Điểm mấu chốt là đạo hàm theo đường cell state: $\partial c_t/\partial c_{t-1}=\mathrm{diag}(f_t)$ (bỏ qua các đường
gián tiếp). Khi $f_t\approx1$, gradient đi qua nhiều bước gần như nguyên vẹn, đúng như số đo ở
Bảng~\ref{tab:ch4-gradflow}. Số tham số là $4H(F+H+1)$, gấp bốn SimpleRNN; với AAPL ($F=6$, $H=64$) bản PyTorch có
\SL{ch4_aapl_torch/lstm/params} tham số (tính cả lớp đầu ra).

\hinhtikz{ch4_lstm}{Ô LSTM: đường cell state $c_t$ (xanh đậm) chỉ chịu phép nhân với cổng quên và phép cộng; ba cổng
sigmoid (đỏ, xanh lá, tím) và ứng viên $\tanh$ (xanh dương).}{ch4-lstm}

% =====================================================================
\section{GRU -- Gated Recurrent Unit}
\label{sec:rnn-gru}

GRU \cite{ch4:cho2014} gộp cell state vào trạng thái ẩn và chỉ dùng hai cổng (Hình~\ref{fig:ch4-gru}):
\begin{align}
  z_t &= \sigma\big(W_z[h_{t-1},x_t]+b_z\big) && \text{cổng cập nhật}\\
  r_t &= \sigma\big(W_r[h_{t-1},x_t]+b_r\big) && \text{cổng đặt lại}\\
  \tilde h_t &= \tanh\big(W_h[r_t\odot h_{t-1},x_t]+b_h\big) && \text{trạng thái ứng viên}\\
  h_t &= (1-z_t)\odot h_{t-1} + z_t\odot\tilde h_t. && \text{nội suy giữa cũ và mới}
\end{align}
Cổng cập nhật $z_t$ đóng vai trò của cả cổng quên và cổng vào của LSTM (ràng buộc tổng bằng 1), cổng đặt lại $r_t$
quyết định trạng thái cũ ảnh hưởng bao nhiêu tới ứng viên. Số tham số $3H(F+H+1)$; với AAPL là
\SL{ch4_aapl_torch/gru/params}.

\hinhtikz{ch4_gru}{Ô GRU với cổng đặt lại $r_t$, cổng cập nhật $z_t$ và phép nội suy
$h_t=(1-z_t)h_{t-1}+z_t\tilde h_t$.}{ch4-gru}

% =====================================================================
\section{So sánh Simple RNN, LSTM và GRU}
\label{sec:rnn-sosanh}

\begin{table}[H]\centering\caption{So sánh ba kiến trúc hồi tiếp cơ bản}\label{tab:ch4-cmp}
\small
\begin{tabular}{L{3.2cm}L{3.5cm}L{3.8cm}L{3.5cm}}
\toprule
Đặc điểm & SimpleRNN & LSTM & GRU \\
\midrule
Bộ nhớ & Trạng thái ẩn $h_t$ & $h_t$ và cell state $c_t$ & Trạng thái ẩn $h_t$ \\
Cổng & Không & Quên, vào, ra & Cập nhật, đặt lại \\
Số tham số lớp & $H(F+H+1)$ & $4H(F+H+1)$ & $3H(F+H+1)$ \\
Phụ thuộc dài hạn & Kém & Tốt & Tốt \\
Chi phí tính & Thấp & Cao nhất & Trung bình \\
Vai trò & Baseline, minh hoạ cơ chế & Bộ nhớ dài, mặc định an toàn & Cân bằng chất lượng/chi phí \\
\bottomrule
\end{tabular}
\end{table}

Chung \& cộng sự \cite{ch4:chung2014} cho thấy LSTM và GRU vượt xa SimpleRNN trên dữ liệu có phụ thuộc dài, nhưng
giữa LSTM và GRU không có kẻ thắng tuyệt đối. Các thực nghiệm dưới đây kiểm tra điều đó trên chuỗi tài chính và kinh
doanh, nơi phụ thuộc dài hạn có thể không tồn tại.

% =====================================================================
\section{Các biến thể và mở rộng quan trọng của RNN}
\label{sec:rnn-bienthe}

\paragraph{RNN hai chiều (Bidirectional RNN).} Hai RNN độc lập đọc chuỗi theo hai chiều; trạng thái mỗi bước là ghép
$[\overrightarrow h_t;\overleftarrow h_t]$ \cite{ch4:schuster1997}. Phù hợp khi đã có toàn bộ chuỗi (gán nhãn, phân
loại văn bản). Trong dự báo, BiLSTM chỉ được phép đọc hai chiều \emph{bên trong cửa sổ quá khứ}, như cách dùng ở
chương này (Hình~\ref{fig:ch4-birnn}).

\hinhtikz{ch4_birnn}{RNN hai chiều: chiều xuôi (xanh) và chiều ngược (hồng) được ghép ở mỗi bước.}{ch4-birnn}

\paragraph{RNN sâu (stacked).} Xếp nhiều tầng hồi tiếp, đầu ra tầng dưới là đầu vào tầng trên (Hình~\ref{fig:ch4-deep});
tầng cao biểu diễn trừu tượng hơn, đổi lại nhiều tham số và dễ quá khớp \cite{ch4:graves2013}.

\hinhtikz{ch4_deep}{RNN sâu ba tầng; mỗi khối 3D là một ô hồi tiếp tại một bước thời gian.}{ch4-deep}

\paragraph{Encoder--decoder và attention.} Bộ mã hoá đọc chuỗi vào thành các trạng thái $e_1..e_n$; bộ giải mã sinh
chuỗi ra. Bản gốc chỉ truyền $e_n$ \cite{ch4:sutskever2014}, tạo nút thắt; attention \cite{ch4:bahdanau2015} cho mỗi
bước giải mã một vector ngữ cảnh riêng $c_j=\sum_i\alpha_{ji}e_i$ với trọng số
$\alpha_{ji}=\mathrm{softmax}_i(\mathrm{score}(s_{j-1},e_i))$ (Hình~\ref{fig:ch4-seq2seq}). Đây là tiền thân trực tiếp
của Transformer \cite{ch4:vaswani2017}.

\hinhtikz{ch4_seq2seq}{Kiến trúc seq2seq với attention: độ dày mũi tên cam tỉ lệ với trọng số chú ý $\alpha$.}{ch4-seq2seq}

\paragraph{Các biến thể khác.} \emph{Stateful RNN} giữ trạng thái giữa các lô liên tiếp; \emph{CNN + RNN} dùng tích
chập 1D trích mẫu cục bộ trước khi RNN mô hình hoá thời gian; \emph{peephole LSTM} cho cổng nhìn trực tiếp $c_{t-1}$;
\emph{mô hình không gian trạng thái} (S4, Mamba \cite{ch4:gu2023}) là hồi tiếp tuyến tính có thể tính song song bằng
phép quét (scan).

% =====================================================================
\section{RNN trong bài toán time-series forecasting}
\label{sec:rnn-forecast}

Dự báo chuỗi thời gian bằng RNN theo quy trình ở Hình~\ref{fig:ch4-pipeline}: nạp dữ liệu, gộp về bước thời gian
mong muốn, khảo sát, chia theo thời gian, chuẩn hoá bằng thống kê của tập train, cắt cửa sổ, huấn luyện, đổi dự báo
về đơn vị gốc, so với baseline và triển khai. Có ba chiến lược dự báo nhiều bước \cite{ch4:hyndman2021}:
\begin{itemize}
  \item \textbf{Một bước (one-step)}: dự báo $t+1$ từ lịch sử tới $t$; đây là chế độ đánh giá chính của chương.
  \item \textbf{Đệ quy (recursive)}: dùng dự báo bước $t+1$ làm đầu vào để dự báo $t+2$, \dots
        (Hình~\ref{fig:ch4-recursive}); một mô hình duy nhất, nhưng sai số tích luỹ theo tầm dự báo. Ứng dụng web dùng
        cách này (mục~\ref{sec:rnn-horizon}).
  \item \textbf{Trực tiếp (direct) hoặc seq2seq}: một mô hình (hoặc một đầu ra) cho mỗi tầm $h$.
\end{itemize}

\hinhtikz{ch4_pipeline}{Quy trình dự báo chuỗi thời gian dùng chung cho ba thực nghiệm.}{ch4-pipeline}
\hinhtikz{ch4_recursive}{Dự báo đệ quy: giá trị vừa dự báo (đỏ) được nối vào cuối cửa sổ để dự báo bước tiếp
theo.}{ch4-recursive}

% =====================================================================
\section{Sliding window / sequence generation}
\label{sec:rnn-window}

Với cửa sổ $w$, mẫu thứ $i$ là $X_i=[x_{i-w},\dots,x_{i-1}]$ và mục tiêu $y_i$ là giá trị tại bước $i$
(Hình~\ref{fig:ch4-window}). $N$ quan sát cho $N-w$ mẫu; cửa sổ trượt bước 1 nên các mẫu kề nhau chồng lên nhau $w-1$
bước, vì vậy \emph{không được xáo trộn rồi chia tập} (cửa sổ test sẽ trùng gần hết với cửa sổ train).

\hinhtikz{ch4_window}{Cửa sổ trượt: mỗi khung màu là một mẫu đầu vào, ô đỏ là giá trị cần dự báo.}{ch4-window}
\maNguon{ch4_data__make_windows}{Tạo cửa sổ trượt many-to-one}

% =====================================================================
\section{Chuẩn hóa dữ liệu}
\label{sec:rnn-chuanhoa}

Các đặc trưng có thang đo rất khác nhau (giá vài trăm USD, khối lượng hàng trăm triệu cổ phiếu), nên phải đưa về cùng
thang trước khi vào mạng. Hai phép phổ biến là Min--Max $x'=(x-x_{\min})/(x_{\max}-x_{\min})$ và z-score
$x'=(x-\mu)/\sigma$. Nguyên tắc bắt buộc: \textbf{thống kê chuẩn hoá chỉ được ước lượng trên tập train}, rồi dùng lại
cho validation và test; fit trên toàn bộ dữ liệu là rò rỉ tương lai.

Với chuỗi giá có xu hướng mạnh, còn một cạm bẫy thứ hai: giá trong tập test có thể vượt xa vùng giá của tập train.
AAPL đóng cửa ở \SL{ch4_eda/aapl/close_first} USD năm 1980 và \SL{ch4_eda/aapl/close_last} USD năm 2024 (gấp
\SL{ch4_eda/aapl/growth_x} lần). Nếu chuẩn hoá Min--Max toàn cục theo train, giá test ánh xạ ra ngoài $[0,1]$, mạng
chưa từng thấy vùng đó và đầu ra ``kẹt trần'' quanh mức cao nhất của train. Đây là nguyên nhân bài mẫu thu được $R^2$
âm cho mọi mô hình trên AAPL. Chương này dùng \textbf{chuẩn hoá tương đối theo cửa sổ} cho dữ liệu giá
(Hình~\ref{fig:ch4-relnorm}): mỗi cột giá được lấy $\log(x_{t-k}/x_t)$ so với bước cuối của chính cửa sổ, và mục tiêu là
log-lợi suất $r_{t+1}=\log(P_{t+1}/P_t)$. Sau khi dự báo $\hat r$, giá được khôi phục bằng
$\hat P_{t+1}=P_t\,e^{\hat r}$. Mục~\ref{sec:rnn-ablation} kiểm chứng định lượng sự khác biệt.

\hinhtikz{ch4_relnorm}{Min--Max toàn cục làm dự báo kẹt ở mức trần của train (trái); chuẩn hoá tương đối đưa mọi cửa
sổ về cùng thang (phải).}{ch4-relnorm}

% =====================================================================
\section{Chia train / validation / test cho chuỗi thời gian}
\label{sec:rnn-split}

Dữ liệu i.i.d.\ có thể chia ngẫu nhiên, nhưng chuỗi thời gian phải chia theo mốc thời gian: train ở trước, validation
ở giữa, test ở cuối (Hình~\ref{fig:ch4-split}), mô phỏng đúng tình huống ``dùng quá khứ dự báo tương lai''. Mọi thực
nghiệm dùng tỉ lệ 70/15/15 trên tập cửa sổ; validation dùng cho dừng sớm và chọn mô hình, test chỉ dùng một lần để
báo cáo.

\hinhtikz{ch4_split}{Chia dữ liệu theo thời gian và nguyên tắc fit bộ chuẩn hoá chỉ trên train.}{ch4-split}
\maNguon{ch4_data__time_split}{Chỉ số cắt train/validation/test theo thứ tự thời gian}

% =====================================================================
\section{Hàm mất mát và tối ưu}
\label{sec:rnn-loss}

Bài toán dự báo là hồi quy, nên dùng sai số bình phương trung bình
\[ \mathcal L_{\text{MSE}}=\frac1n\sum_{i=1}^n\big(y_i-\hat y_i\big)^2 \]
trên mục tiêu đã chuẩn hoá. Tối ưu bằng Adam \cite{ch4:kingma2015}, kết hợp động lượng và tốc độ học thích nghi theo
từng tham số:
\begin{align*}
  m_t&=\beta_1 m_{t-1}+(1-\beta_1)g_t,\quad v_t=\beta_2v_{t-1}+(1-\beta_2)g_t^2,\\
  \hat m_t&=\frac{m_t}{1-\beta_1^t},\quad \hat v_t=\frac{v_t}{1-\beta_2^t},\quad
  \theta_t=\theta_{t-1}-\eta\,\frac{\hat m_t}{\sqrt{\hat v_t}+\epsilon},
\end{align*}
với $\eta=10^{-3}$, $\beta_1=0{,}9$, $\beta_2=0{,}999$. Tốc độ học được giảm một nửa khi validation loss không cải
thiện sau 6 epoch (ReduceLROnPlateau).

% =====================================================================
\section{Đánh giá mô hình forecasting}
\label{sec:rnn-metrics}

\begin{table}[H]\centering\caption{Các chỉ số đánh giá dự báo (tính trên đơn vị gốc)}\label{tab:ch4-metric}
\small
\begin{tabular}{L{2.6cm}L{5.4cm}L{6.2cm}}
\toprule
Chỉ số & Công thức & Cách đọc \\
\midrule
MAE & $\frac1n\sum|y_i-\hat y_i|$ & Sai số trung bình, cùng đơn vị mục tiêu; bền với ngoại lai \\[4pt]
RMSE & $\sqrt{\frac1n\sum(y_i-\hat y_i)^2}$ & Phạt nặng sai số lớn; cùng đơn vị \\[4pt]
MAPE & $\frac{100}{n}\sum\big|\frac{y_i-\hat y_i}{y_i}\big|$ & Phần trăm; không ổn định khi $y$ gần 0 \\[4pt]
$R^2$ & $1-\frac{\sum(y_i-\hat y_i)^2}{\sum(y_i-\bar y)^2}$ & 1 là hoàn hảo; âm nghĩa là kém hơn đoán trung bình \\[4pt]
Đúng hướng & $\frac1n\sum\mathbf{1}\big[\mathrm{sgn}(\hat y_i-y_{i-1})=\mathrm{sgn}(y_i-y_{i-1})\big]$ & Tỉ lệ đoán đúng chiều tăng/giảm; 50\% là ngẫu nhiên \\[4pt]
RMSE / naive & $\mathrm{RMSE}_{\text{mô hình}}/\mathrm{RMSE}_{\hat y_i=y_{i-1}}$ & $<1$ mới thực sự có kỹ năng dự báo \\
\bottomrule
\end{tabular}
\end{table}

\begin{luuy}{Vì sao phải so với dự báo ngây thơ}
Với giá tài sản, $R^2$ trên mức giá gần như luôn rất cao (giá hôm nay đã giải thích gần hết giá ngày mai), nên $R^2$
cao \emph{không} chứng minh mô hình giỏi. Thước đo trung thực là so với \textbf{dự báo ngây thơ} (naive, persistence)
$\hat y_{t+1}=y_t$; mô hình chỉ có giá trị nếu RMSE của nó nhỏ hơn RMSE naive. Cách đọc này bám theo lời khuyên ``đọc
biểu đồ dự báo, không chỉ đọc bảng'' của các bài tham khảo: một đường dự báo bám sát thực tế nhưng trễ một phiên chính
là dự báo naive.
\end{luuy}

\maNguon{ch4_metrics__forecast_metrics}{Tính MAE, RMSE, MAPE, $R^2$ và độ chính xác hướng}

% =====================================================================
\section{Overfitting, regularization và early stopping}
\label{sec:rnn-overfit}

Quá khớp xảy ra khi mô hình học thuộc nhiễu của train. Với RNN trên chuỗi ngắn (e-commerce chỉ có
\SL{ch4_eda/ecom/days} ngày), nguy cơ rất cao. Các biện pháp dùng trong chương: Dropout 0,2 trước lớp đầu ra
\cite{ch4:srivastava2014}; dừng sớm (early stopping) với kiên nhẫn 15 epoch theo validation loss và khôi phục trọng số
tốt nhất (Hình~\ref{fig:ch4-earlystop}); giảm tốc độ học khi chững lại; cắt chuẩn gradient. Các lựa chọn khác gồm
recurrent dropout, weight decay và tăng dữ liệu bằng nhiều chuỗi.

\hinhtikz{ch4_earlystop}{Dừng sớm: validation loss đạt cực tiểu rồi tăng khi mô hình bắt đầu quá khớp; trọng số tại
epoch tốt nhất được khôi phục.}{ch4-earlystop}

% =====================================================================
\section{Hyperparameter quan trọng của RNN}
\label{sec:rnn-sieuthamso}

\begin{table}[H]\centering\caption{Siêu tham số chính và giá trị dùng trong chương}\label{tab:ch4-hparams}
\small
\begin{tabular}{L{3.3cm}L{6.7cm}L{4.2cm}}
\toprule
Siêu tham số & Ảnh hưởng & Giá trị dùng \\
\midrule
sequence\_length & Độ dài lịch sử mô hình nhìn thấy (mục~\ref{sec:rnn-ablation} khảo sát 10--90) & 30 \\
hidden\_size & Năng lực biểu diễn; lớn quá tốn tài nguyên, quá khớp & 64 \\
num\_layers & Độ sâu mạng & 1 \\
dropout & Chống quá khớp & 0,2 \\
learning\_rate & Bước cập nhật & $10^{-3}$, giảm ×0,5 khi chững \\
batch\_size & Tốc độ và độ nhiễu của gradient & 64 \\
epochs / patience & Số lượt quét tối đa / số epoch chờ trước khi dừng & 150 / 15 \\
clipnorm & Ngưỡng cắt chuẩn gradient & 1,0 \\
bidirectional & Đọc hai chiều trong cửa sổ & chỉ ở BiLSTM \\
optimizer, loss & Cách cập nhật, mục tiêu tối ưu & Adam, MSE \\
\bottomrule
\end{tabular}
\end{table}

% =====================================================================
\section{Cài đặt RNN từ đầu bằng NumPy}
\label{sec:rnn-scratch}

Trước khi dùng thư viện, ta cài đặt một SimpleRNN many-to-one chỉ bằng NumPy để thấy rõ từng phép tính của
mục~\ref{sec:rnn-cautruc} và~\ref{sec:rnn-bptt}. Lớp \texttt{RNNScratch} giữ năm tham số $W_{xh}, W_{hh}, b_h, W_{hy},
b_y$, khởi tạo theo phân phối chuẩn với độ lệch $1/\sqrt{H}$.

\maNguon{ch4_scratch__RNNScratch.forward}{Lan truyền xuôi: lặp qua các bước, lưu mọi $h_t$ để dùng cho BPTT}

Lan truyền ngược cài đặt đúng công thức~\eqref{eq:bptt-delta}--\eqref{eq:bptt-w}: tín hiệu lỗi $\delta$ đi từ bước cuối
về bước đầu, mỗi bước nhân với $\mathrm{diag}(1-h_t^2)$ rồi với $W_{hh}^\top$; gradient của $W_{hh}$, $W_{xh}$ được cộng
dồn. Hàm còn ghi lại $\|\delta_t\|$ ở từng bước để quan sát gradient biến mất ngay trên mô hình thật.

\maNguon{ch4_scratch__RNNScratch.backward}{BPTT viết tay: cộng dồn gradient theo thời gian và ghi chuẩn $\|\delta_t\|$}

\subsection{Kiểm tra gradient bằng sai phân số}

Một BPTT viết tay rất dễ sai chỉ số. Cách kiểm tra chuẩn là so gradient giải tích với sai phân trung tâm
\begin{gather*}
  \frac{\partial L}{\partial\theta_j}\approx\frac{L(\theta+\epsilon e_j)-L(\theta-\epsilon e_j)}{2\epsilon},\\
  \text{sai số tương đối}=\frac{|g_{\text{giải tích}}-g_{\text{số}}|}{\max(|g_{\text{giải tích}}|+|g_{\text{số}}|,\,10^{-12})}.
\end{gather*}
Hàm \texttt{grad\_check} lấy ngẫu nhiên 12 phần tử của mỗi tham số, $\epsilon=10^{-5}$.

\maNguon{ch4_scratch__grad_check}{Kiểm tra gradient BPTT bằng sai phân trung tâm}
\bang{ch4_gradcheck}{Sai số tương đối lớn nhất giữa gradient BPTT và sai phân số}{ch4-gradcheck}

Sai số lớn nhất chỉ \SL{ch4_scratch/max_rel_err}, nằm trong vùng sai số làm tròn của số thực 64 bit: BPTT viết tay
đúng. Để chắc chắn công thức xuôi cũng khớp với thư viện, trọng số NumPy được chép sang \texttt{torch.nn.RNN} và hai
đầu ra được so trực tiếp; chênh lệch lớn nhất là \SL{ch4_scratch/torch_diff}.

\maNguon{ch4_scratch_run__torch_equivalence}{So khớp RNN NumPy với \texttt{torch.nn.RNN} khi dùng chung trọng số}

\subsection{Huấn luyện trên chuỗi sin có nhiễu}

Dữ liệu thử là chuỗi $\sin$ cộng nhiễu Gauss, cửa sổ \SL{ch4_scratch/window} bước, \SL{ch4_scratch/n_train} mẫu
train, \SL{ch4_scratch/n_val} mẫu validation và \SL{ch4_scratch/n_test} mẫu test. Mô hình có \SL{ch4_scratch/hidden}
đơn vị ẩn, tổng cộng \SL{ch4_scratch/params} tham số, huấn luyện bằng SGD mini-batch có cắt chuẩn gradient trong
\SL{ch4_scratch/epochs} epoch.

\maNguon{ch4_scratch_run__sine_series}{Sinh chuỗi sin có nhiễu cho thí nghiệm RNN từ đầu}
\maNguon{ch4_scratch__sgd_train}{Vòng huấn luyện SGD mini-batch có cắt chuẩn gradient cho \texttt{RNNScratch}}

\hinhpdf{ch4_scratch_train}{1}{RNN NumPy: MSE train/validation theo epoch (trái) và dự báo một bước trên test
(phải).}{ch4-scratch-train}

\nhanxet{MSE cuối cùng là \SL{ch4_scratch/final_train} trên train và \SL{ch4_scratch/final_val} trên validation, hai
đường gần nhau nên mô hình không quá khớp. RMSE trên test là \SL{ch4_scratch/rmse}, nhỏ hơn dự báo naive
(\SL{ch4_scratch/rmse_naive}) gần một nửa: khác với chuỗi tài chính ở các mục sau, chuỗi sin có cấu trúc tuần hoàn
thật mà RNN học được.}

\hinhpdf{ch4_scratch_dh}{0.92}{Chuẩn $\|\delta_t\|=\|\partial L/\partial h_t\|$ của RNN NumPy đã huấn luyện, đi ngược từ
bước cuối về bước đầu của cửa sổ.}{ch4-scratch-dh}

\nhanxet{Ngay cả trên cửa sổ chỉ \SL{ch4_scratch/window} bước, chuẩn gradient tại bước cuối cửa sổ ($T-t=0$) là
\SL{ch4_scratch/dh_first}, còn khi lùi về bước đầu cửa sổ chỉ còn \SL{ch4_scratch/dh_last}: tín hiệu lỗi suy giảm hơn
một bậc. Đây là bằng chứng trực tiếp cho phân tích ở mục~\ref{sec:rnn-vanish}.}

% =====================================================================
\section{Thiết lập chung cho ba thực nghiệm}
\label{sec:rnn-setup}

\subsection{Dữ liệu}

Ba bộ dữ liệu thực, tải từ Kaggle bằng \texttt{kagglehub} (Bảng~\ref{tab:ch4-datasets}):
\begin{itemize}
  \item \textbf{E-commerce Customer Behavior \& Sales -- TR} \cite{ch4:data_ecom}: \SL{ch4_eda/ecom/orders} đơn hàng của
        một sàn thương mại điện tử Thổ Nhĩ Kỳ, \SL{ch4_eda/ecom/raw_cols} cột. Đơn được gộp thành chuỗi ngày.
  \item \textbf{AAPL Stock Performance} \cite{ch4:data_aapl}: \SL{ch4_eda/aapl/n_obs} phiên giao dịch cổ phiếu Apple từ
        \SL{ch4_eda/aapl/start} đến \SL{ch4_eda/aapl/end}.
  \item \textbf{Gold Price Prediction Dataset} \cite{ch4:data_gold}: \SL{ch4_eda/gold/n_obs} phiên của quỹ vàng GLD
        cùng \SL{ch4_eda/gold/cols_raw} cột thị trường liên quan; dùng \SL{ch4_eda/gold/cols_used} cột.
\end{itemize}

\bang{ch4_datasets}{Ba bộ dữ liệu chuỗi thời gian của chương}{ch4-datasets}
\maNguon{ch4_data__download_all}{Tải ba bộ dữ liệu Kaggle vào \texttt{data/ch4}}
\maNguon{ch4_data__experiment}{Đóng gói mỗi thực nghiệm: khung dữ liệu, cột mục tiêu, kiểu chuẩn hoá, đơn vị}

\subsection{Tiền xử lý và cửa sổ}

Lớp \texttt{Prepared} thực hiện toàn bộ đường ống ở Hình~\ref{fig:ch4-pipeline}: cắt cửa sổ 30 bước, chia 70/15/15
theo thời gian trên tập cửa sổ, ước lượng thống kê chuẩn hoá chỉ trên phần train, và cung cấp các baseline naive,
trung bình trượt để so sánh. Có ba chế độ chuẩn hoá: \texttt{level} (z-score trên mức, dùng cho e-commerce vốn không
có xu hướng), \texttt{relative} (chuẩn hoá tương đối theo cửa sổ, dùng cho giá) và \texttt{minmax} (Min--Max toàn
cục, chỉ để làm thí nghiệm đối chứng ở mục~\ref{sec:rnn-ablation}).

\maDai{ch4_data__Prepared.__init__}{Khởi tạo \texttt{Prepared}: cửa sổ, chia tập theo thời gian và chọn kiểu chuẩn hoá}
\maNguon{ch4_data__Prepared._level}{Chuẩn hoá z-score trên mức, thống kê lấy từ train}
\maNguon{ch4_data__Prepared._relative}{Chuẩn hoá tương đối theo cửa sổ: log-tỉ số so với bước cuối, mục tiêu là
log-lợi suất}
\maNguon{ch4_data__Prepared.to_original}{Đổi dự báo đã chuẩn hoá về đơn vị gốc}
\maNguon{ch4_data__Prepared.naive}{Dự báo ngây thơ: giá trị của bước cuối cửa sổ}
\maNguon{ch4_data__Prepared.moving_average}{Baseline trung bình trượt $k$ ngày}

\bang{ch4_shapes}{Kích thước tensor đầu vào sau khi cắt cửa sổ}{ch4-shapes}

\subsection{Mô hình}

Cả hai framework dùng cùng một kiến trúc: một lớp hồi tiếp 64 đơn vị (SimpleRNN, LSTM, GRU, hoặc BiLSTM $2\times64$),
lấy trạng thái ẩn ở bước cuối, Dropout 0,2 và một lớp tuyến tính ra một giá trị. Bản PyTorch \cite{ch4:paszke2019}
chạy trên GPU, bản Keras \cite{ch4:chollet2021} chạy trên CPU sáu luồng.

\maNguon{ch4_torch__SeqRegressor}{Mô hình hồi quy chuỗi trong PyTorch, chọn kiến trúc bằng tham số \texttt{arch}}
\maNguon{ch4_torch__train_epoch}{Một epoch huấn luyện PyTorch có cắt chuẩn gradient}
\maNguon{ch4_torch__predict}{Suy luận theo lô trong PyTorch}
\clearpage
\maDai{ch4_torch__fit}{Vòng huấn luyện PyTorch: Adam, ReduceLROnPlateau, dừng sớm và khôi phục trọng số tốt nhất}
\maNguon{ch4_keras__build_keras_model}{Cùng kiến trúc trong Keras: \texttt{SimpleRNN}/\texttt{LSTM}/\texttt{GRU}/\texttt{Bidirectional}}
\maNguon{ch4_keras__fit_keras}{Huấn luyện Keras với \texttt{EarlyStopping}, khôi phục trọng số tốt nhất}

\bang{ch4_params}{Số tham số và tốc độ huấn luyện của bốn kiến trúc}{ch4-params}

\subsection{Quy trình đánh giá}

Mỗi mô hình được huấn luyện tối đa 150 epoch, dừng sớm theo validation loss, rồi dự báo một bước trên test. Dự báo
được đổi về đơn vị gốc trước khi tính MAE, RMSE, MAPE, $R^2$ và độ chính xác hướng; tất cả được so với dự báo naive
trên cùng tập test.

\maNguon{ch4_train__evaluate}{Đổi dự báo về đơn vị gốc và tính chỉ số trên validation/test}
\maDai{ch4_train__run}{Kịch bản huấn luyện: 3 thực nghiệm × 4 kiến trúc cho một framework, ghi JSON, NPZ và mô hình}

% =====================================================================
\section{Thực nghiệm 1: dự báo doanh thu ngày của sàn thương mại điện tử}
\label{sec:rnn-tn1}

\subsection{Bài toán}
Dự báo tổng doanh thu \texttt{Total\_Amount} (TRY) của ngày kế tiếp từ 30 ngày gần nhất của chín chuỗi hành vi: doanh
thu, số giao dịch, số lượng, đơn giá trung bình, chiết khấu, thời lượng phiên, số trang xem, điểm đánh giá và tỉ lệ
khách quay lại. Đây là bài toán dự báo nhu cầu kinh điển cho kế hoạch tồn kho và nhân sự.

\hinhtikz{ch4_aggregate}{Gộp các đơn hàng lẻ thành chuỗi ngày: mỗi ngày là một vector chín đặc trưng (tổng, đếm,
trung bình).}{ch4-aggregate}
\maNguon{ch4_data__ecommerce_daily}{Gộp đơn hàng thành chín chuỗi theo ngày}

\subsection{Khám phá dữ liệu}

Dữ liệu sạch: \SL{ch4_eda/ecom/missing} giá trị thiếu, \SL{ch4_eda/ecom/dup} dòng trùng, không có ngày nào trống
đơn. Gộp theo ngày cho \SL{ch4_eda/ecom/days} ngày từ \SL{ch4_eda/ecom/start} đến \SL{ch4_eda/ecom/end}, trung bình
\SL{ch4_eda/ecom/orders_per_day} đơn mỗi ngày. Doanh thu ngày trung bình \SL{ch4_eda/ecom/target/mean} TRY, độ lệch
chuẩn \SL{ch4_eda/ecom/target/std}, dao động từ \SL{ch4_eda/ecom/target/min} đến \SL{ch4_eda/ecom/target/max}.

\hinhpdf{ch4_eda_ecom_series}{1}{Các chuỗi hành vi theo ngày của sàn thương mại điện tử.}{ch4-eda-ecom-series}
\hinhpdf{ch4_eda_ecom_hist}{1}{Phân phối của chín đặc trưng ngày; doanh thu lệch phải mạnh.}{ch4-eda-ecom-hist}
\hinhpdf{ch4_eda_ecom_corr}{0.82}{Ma trận tương quan giữa các đặc trưng ngày.}{ch4-eda-ecom-corr}

\nhanxet{Doanh thu tương quan dương với số lượng (\SL{ch4_eda/ecom/corr_amount_qty}) và số giao dịch
(\SL{ch4_eda/ecom/corr_amount_tx}), nhưng gần như không liên quan tới điểm đánh giá
(\SL{ch4_eda/ecom/corr_amount_rating}). Như vậy cấu trúc đáng khai thác nằm ở quy mô đơn hàng trong ngày, không phải ở
các chỉ số trải nghiệm.}

\hinhpdf{ch4_eda_ecom_season}{1}{Doanh thu trung bình theo thứ, tổng theo tháng và hàm tự tương quan (ACF) của doanh
thu ngày.}{ch4-eda-ecom-season}

\begin{phathien}{Phát hiện: doanh thu gần như không tự tương quan}
Hệ số tự tương quan bậc 1 chỉ \SL{ch4_eda/ecom/acf1}, bậc 7 là \SL{ch4_eda/ecom/acf7}: doanh thu hôm nay hầu như
không cho biết gì về doanh thu ngày mai. Chênh lệch theo thứ có tồn tại (ngày cao nhất \SL{ch4_eda/ecom/dow_max} so với
thấp nhất \SL{ch4_eda/ecom/dow_min} gấp \SL{ch4_eda/ecom/dow_ratio} lần), nhưng nhỏ so với nhiễu. Hệ quả: dự báo naive
``ngày mai bằng hôm nay'' sẽ rất tệ, vì nó lặp lại nguyên nhiễu của hôm nay; một mô hình chỉ cần học được \emph{mức
trung bình} đã thắng naive.
\end{phathien}

\maNguon{ch4_eda__acf}{Hàm tự tương quan dùng trong phần khám phá dữ liệu}
\hinhpdf{ch4_eda_ecom_mix}{1}{Cơ cấu doanh thu theo danh mục (nhiều nhất là \SL{ch4_eda/ecom/top_category}) và đơn
theo thiết bị.}{ch4-eda-ecom-mix}
\hinhpdf{ch4_eda_ecom_split}{1}{Chia train/validation/test theo thời gian: train tới \SL{ch4_eda/ecom/split/train_end},
validation tới \SL{ch4_eda/ecom/split/val_end}, test từ \SL{ch4_eda/ecom/split/test_start}.}{ch4-eda-ecom-split}

\subsection{Huấn luyện}

\hinhpdf{ch4_ecommerce_loss_torch}{1}{E-commerce: MSE train và validation theo epoch của bốn kiến trúc PyTorch; chấm
đỏ là epoch có validation loss nhỏ nhất.}{ch4-ecom-loss}
\hinhpdf{ch4_ecommerce_loss_keras}{1}{E-commerce: đường học của bốn kiến trúc Keras.}{ch4-ecom-loss-keras}

\nhanxet{Với chỉ \SL{ch4_eda/shapes/ecommerce/train/0} cửa sổ train, validation loss chạm đáy sau vài chục epoch rồi
đi ngang hoặc tăng trong khi train loss vẫn giảm: dấu hiệu quá khớp điển hình mà dừng sớm xử lý. LSTM PyTorch dừng ở
epoch \SL{ch4_ecommerce_torch/lstm/epochs}, GRU Keras ở epoch \SL{ch4_ecommerce_keras/gru/epochs}.}

\subsection{Kết quả}

\bang{ch4_ecommerce_results}{E-commerce: kết quả dự báo doanh thu ngày trên tập test}{ch4-ecom-results}

\hinhpdf{ch4_ecommerce_pred_torch}{1}{E-commerce: dự báo một bước của bốn kiến trúc PyTorch trên toàn bộ tập
test.}{ch4-ecom-pred-torch}
\hinhpdf{ch4_ecommerce_pred_keras}{1}{E-commerce: dự báo một bước của bốn kiến trúc Keras trên toàn bộ tập
test.}{ch4-ecom-pred-keras}
\hinhpdf{ch4_ecommerce_zoom}{1}{E-commerce: 30 ngày cuối của test, so LSTM/GRU của hai framework với naive và trung
bình trượt 7 ngày.}{ch4-ecom-zoom}

\nhanxet{Naive có RMSE \SL{ch4_ecommerce_torch/naive/test/rmse} TRY, trung bình trượt 7 ngày
\SL{ch4_summary/ecommerce/ma7/rmse} và 30 ngày \SL{ch4_summary/ecommerce/ma30/rmse}. Cả \SL{ch4_summary/ecommerce/n_beat_naive}
mô hình RNN đều thắng naive; tốt nhất là LSTM PyTorch với RMSE \SL{ch4_ecommerce_torch/lstm/test/rmse}
($R^2$ = \SL{ch4_ecommerce_torch/lstm/test/r2}), Keras tốt nhất là GRU với \SL{ch4_summary/ecommerce/keras_best_rmse}.
Quan trọng hơn, RNN còn thắng cả trung bình trượt 30 ngày, nghĩa là nó không chỉ ``làm trơn'' mà dùng được thông tin
của các đặc trưng phụ (số giao dịch, số lượng). Hình~\ref{fig:ch4-ecom-zoom} cho thấy đường dự báo bám mức trung bình
và bỏ qua phần lớn dao động ngày: đó là hành vi đúng khi ACF gần 0. $R^2$ chỉ khoảng 0,1 nhắc rằng phần lớn biến động
doanh thu ngày là ngẫu nhiên và không mô hình nào dự báo được.}

\hinhpdf{ch4_ecommerce_resid}{1}{E-commerce: phần dư của mô hình tốt nhất theo thời gian và phân phối của nó.}{ch4-ecom-resid}
\hinhpdf{ch4_ecommerce_bars}{1}{E-commerce: RMSE và MAE của bốn kiến trúc trên hai framework, đường đứt là naive,
đường chấm là trung bình trượt 7 ngày.}{ch4-ecom-bars}

\nhanxet{Phần dư dao động quanh 0, không có xu hướng theo thời gian, phân phối lệch phải do những ngày doanh thu đột
biến mà mô hình (hợp lý) không đoán trước được. Độ chính xác hướng của LSTM PyTorch là
\SL{ch4_ecommerce_torch/lstm/test/dir_accPct}\%, cao vì doanh thu có tính hồi quy về trung bình: sau một ngày cao,
dự báo ``giảm'' thường đúng.}

% =====================================================================
\section{Thực nghiệm 2: dự báo giá cổ phiếu AAPL}
\label{sec:rnn-tn2}

\subsection{Bài toán}
Dự báo giá đóng cửa \texttt{close} của phiên kế tiếp từ 30 phiên gần nhất với sáu đặc trưng open, high, low, close,
adjclose, volume. Đây là bài toán được nhiều bài tham khảo chọn; các bài này thường báo cáo $R^2$ âm hoặc đường dự báo
lệch hẳn khỏi giá thật ở cuối chuỗi, và mục này giải thích vì sao.

\maNguon{ch4_data__aapl_frame}{Đọc dữ liệu AAPL, sắp xếp theo ngày, giữ sáu cột số}

\subsection{Khám phá dữ liệu}

Giá đóng cửa đi từ \SL{ch4_eda/aapl/close_first} USD lên \SL{ch4_eda/aapl/close_last} USD (đỉnh
\SL{ch4_eda/aapl/close_max}), tức gấp \SL{ch4_eda/aapl/growth_x} lần trong 44 năm. Lợi suất ngày trung bình
\SL{ch4_eda/aapl/ret_mean_pct}\%, độ lệch chuẩn \SL{ch4_eda/aapl/ret_sd_pct}\%, nhưng có đuôi rất dày: độ nhọn dư
\SL{ch4_eda/aapl/ret_kurt}, phiên giảm mạnh nhất \SL{ch4_eda/aapl/ret_min_pct}\% ngày
\SL{ch4_eda/aapl/ret_min_date}.

\hinhpdf{ch4_eda_aapl_overview}{1}{AAPL: giá đóng cửa (thang log), chênh lệch close--open, biên độ high--low và khối
lượng.}{ch4-eda-aapl-overview}
\hinhpdf{ch4_eda_aapl_hist}{1}{Phân phối của sáu đặc trưng AAPL; khối lượng vẽ theo $\log(1+x)$.}{ch4-eda-aapl-hist}
\hinhpdf{ch4_eda_aapl_corr}{0.75}{Tương quan giữa các đặc trưng AAPL: bốn cột giá gần như trùng nhau.}{ch4-eda-aapl-corr}
\hinhpdf{ch4_eda_aapl_yearly}{1}{AAPL: giá cuối năm và lợi suất từng năm; năm tốt nhất \SL{ch4_eda/aapl/best_yearRaw}
(\SL{ch4_eda/aapl/best_year_pct}\%), tệ nhất \SL{ch4_eda/aapl/worst_yearRaw} (\SL{ch4_eda/aapl/worst_year_pct}\%).}{ch4-eda-aapl-yearly}
\hinhpdf{ch4_eda_aapl_returns}{1}{AAPL: biến động năm hoá trên cửa sổ 60 phiên và đuôi dày của lợi suất ngày so với
phân phối chuẩn.}{ch4-eda-aapl-returns}

\begin{phathien}{Phát hiện: lợi suất không tự tương quan, độ biến động thì có}
ACF bậc 1 của lợi suất chỉ \SL{ch4_eda/aapl/acf1_ret}, trong khi ACF bậc 1 của \emph{trị tuyệt đối} lợi suất là
\SL{ch4_eda/aapl/acf1_abs}. Hướng đi của giá gần như không dự báo được từ quá khứ, phù hợp với giả thuyết thị trường
hiệu quả \cite{ch4:fama1970}; còn độ lớn dao động thì có tính tụ cụm (biến động cao đi liền biến động cao). Ta nên kỳ
vọng RNN không thắng naive đáng kể về giá.
\end{phathien}

\hinhpdf{ch4_eda_aapl_split}{1}{AAPL: chia theo thời gian (thang log). Giá test (từ \SL{ch4_eda/aapl/split/test_start})
cao hơn hẳn mọi giá trong train.}{ch4-eda-aapl-split}

\subsection{Huấn luyện}

\hinhpdf{ch4_aapl_loss_torch}{1}{AAPL: MSE train và validation trên log-lợi suất đã chuẩn hoá (PyTorch).}{ch4-aapl-loss}
\hinhpdf{ch4_aapl_loss_keras}{1}{AAPL: đường học của bốn kiến trúc Keras.}{ch4-aapl-loss-keras}

\nhanxet{Validation loss gần như đi ngang quanh 1 ngay từ đầu (phương sai của lợi suất đã chuẩn hoá): mạng không tìm
được tín hiệu nào tốt hơn ``dự báo lợi suất bằng 0'' và dừng sớm sau vài chục epoch. Đây là đường học đặc trưng của
một chuỗi gần với bước ngẫu nhiên.}

\subsection{Kết quả}

\bang{ch4_aapl_results}{AAPL: kết quả dự báo giá đóng cửa trên tập test}{ch4-aapl-results}
\hinhpdf{ch4_aapl_pred_torch}{1}{AAPL: dự báo một bước của bốn kiến trúc PyTorch trên test
(\SL{ch4_eda/aapl/split/test_start} -- \SL{ch4_eda/aapl/end}).}{ch4-aapl-pred-torch}
\hinhpdf{ch4_aapl_pred_keras}{1}{AAPL: dự báo một bước của bốn kiến trúc Keras trên test.}{ch4-aapl-pred-keras}
\hinhpdf{ch4_aapl_zoom}{1}{AAPL: 60 phiên cuối của test.}{ch4-aapl-zoom}

\nhanxet{Mọi mô hình có $R^2$ khoảng \SL{ch4_aapl_torch/lstm/test/r2} và đường dự báo trùng khít đường thực tế ở
Hình~\ref{fig:ch4-aapl-pred-torch}. Nhưng naive cũng có RMSE \SL{ch4_aapl_torch/naive/test/rmse} USD, còn mô hình tốt
nhất đạt \SL{ch4_summary/aapl/best/rmse}, tỉ số chỉ \SL{ch4_summary/aapl/best_ratio}. Độ chính xác hướng của mô hình
tốt nhất là \SL{ch4_summary/aapl/best/dir_accPct}\%, tức tương đương tung đồng xu. Nhìn kỹ ở
Hình~\ref{fig:ch4-aapl-zoom}, đường dự báo chính là đường thực tế dịch sang phải một phiên. Kết luận trung thực: RNN học
được đúng điều duy nhất có thể học từ giá quá khứ, ``giá ngày mai xấp xỉ giá hôm nay'', và không học được hướng.
Ngược lại với bài mẫu, mô hình không bị kẹt ở trần của train nhờ chuẩn hoá tương đối; mục~\ref{sec:rnn-ablation} đo
định lượng khác biệt này.}

\hinhpdf{ch4_aapl_resid}{1}{AAPL: phần dư theo thời gian và phân phối; độ lớn phần dư tăng theo mức giá và theo các
giai đoạn biến động.}{ch4-aapl-resid}
\hinhpdf{ch4_aapl_bars}{1}{AAPL: RMSE và MAE theo kiến trúc và framework; đường đứt là naive.}{ch4-aapl-bars}

% =====================================================================
\section{Thực nghiệm 3: dự báo giá quỹ vàng GLD}
\label{sec:rnn-tn3}

\subsection{Bài toán}
Dự báo giá đóng cửa của quỹ vàng GLD ở phiên kế tiếp. Khác với AAPL, mô hình được cho thêm thông tin \emph{liên thị
trường}: chỉ số S\&P~500, Dow Jones, tỉ giá EUR/USD, dầu Brent, chỉ số USD, hợp đồng bạc, quỹ cổ phiếu mỏ vàng GDX và
quỹ dầu USO, tổng cộng \SL{ch4_eda/gold/cols_used} đặc trưng. Câu hỏi: thông tin bổ sung có giúp RNN vượt naive không?

\maNguon{ch4_data__gold_frame}{Chọn 13 cột thị trường từ bộ dữ liệu vàng}

\subsection{Khám phá dữ liệu}

GLD dao động từ \SL{ch4_eda/gold/target/min} đến \SL{ch4_eda/gold/target/max} USD trong giai đoạn
\SL{ch4_eda/gold/start} -- \SL{ch4_eda/gold/end}, sụt giảm lớn nhất \SL{ch4_eda/gold/drawdown_pct}\% sau đỉnh 2012;
độ lệch chuẩn lợi suất ngày \SL{ch4_eda/gold/ret_sd_pct}\%, thấp hơn nhiều so với AAPL.

\hinhpdf{ch4_eda_gold_series}{1}{Vàng GLD và các thị trường liên quan 2011--2018, chuẩn hoá về 100 ở phiên đầu.}{ch4-eda-gold-series}
\hinhpdf{ch4_eda_gold_corr}{0.85}{Tương quan lợi suất ngày giữa GLD và 12 chuỗi còn lại.}{ch4-eda-gold-corr}

\begin{phathien}{Phát hiện: tương quan cùng phiên mạnh, nhưng không phải dự báo}
Lợi suất GLD tương quan mạnh với GDX (\SL{ch4_eda/gold/corr_ret/GDX}) và bạc (\SL{ch4_eda/gold/corr_ret/Silver_Fut}),
ngược chiều với chỉ số USD (\SL{ch4_eda/gold/corr_ret/USD_Index}), gần như độc lập với S\&P~500
(\SL{ch4_eda/gold/corr_ret/SP500}). Nhưng đó là tương quan \emph{cùng phiên}: để biết GDX hôm nay thì cũng phải đợi tới
cuối phiên, khi giá vàng đã biết. Thông tin có ích cho dự báo ngày mai phải là tương quan \emph{trễ}, và nó yếu hơn
nhiều.
\end{phathien}

\hinhpdf{ch4_eda_gold_rollcorr}{1}{Tương quan trượt 120 phiên giữa lợi suất GLD và các thị trường: quan hệ thay đổi
theo thời gian.}{ch4-eda-gold-rollcorr}
\hinhpdf{ch4_eda_gold_hist}{1}{Phân phối của các đặc trưng thị trường vàng.}{ch4-eda-gold-hist}
\hinhpdf{ch4_eda_gold_split}{1}{GLD: chia theo thời gian; test bắt đầu \SL{ch4_eda/gold/split/test_start}.}{ch4-eda-gold-split}

\subsection{Huấn luyện}

\hinhpdf{ch4_gold_loss_torch}{1}{Vàng GLD: MSE train và validation theo epoch (PyTorch).}{ch4-gold-loss}
\hinhpdf{ch4_gold_loss_keras}{1}{Vàng GLD: đường học của bốn kiến trúc Keras.}{ch4-gold-loss-keras}

\nhanxet{Train loss giảm rõ (13 đặc trưng cho mạng nhiều cách khớp nhiễu của train), nhưng validation loss gần như
không giảm: khoảng cách train--validation lớn là dấu hiệu của quá khớp, và dừng sớm giữ lại trọng số của những epoch
đầu.}

\subsection{Kết quả}

\bang{ch4_gold_results}{Vàng GLD: kết quả dự báo giá đóng cửa trên tập test}{ch4-gold-results}
\hinhpdf{ch4_gold_pred_torch}{1}{Vàng GLD: dự báo một bước của bốn kiến trúc PyTorch trên test.}{ch4-gold-pred-torch}
\hinhpdf{ch4_gold_pred_keras}{1}{Vàng GLD: dự báo một bước của bốn kiến trúc Keras trên test.}{ch4-gold-pred-keras}
\hinhpdf{ch4_gold_zoom}{1}{Vàng GLD: 60 phiên cuối của test.}{ch4-gold-zoom}

\nhanxet{Naive đạt RMSE \SL{ch4_gold_torch/naive/test/rmse} USD; LSTM PyTorch \SL{ch4_gold_torch/lstm/test/rmse} và
LSTM Keras \SL{ch4_gold_keras/lstm/test/rmse}. Chỉ \SL{ch4_summary/gold/n_beat_naive} trên 8 mô hình nhỉnh hơn naive và
chênh lệch dưới 1\%. $R^2$ khoảng \SL{ch4_gold_torch/lstm/test/r2} nghe ấn tượng nhưng, như với AAPL, gần như toàn bộ
đến từ thông tin ``giá hôm nay''. Thêm 12 chuỗi liên thị trường không đem lại kỹ năng dự báo một bước, đúng như phân
tích tương quan trễ ở trên.}

\hinhpdf{ch4_gold_resid}{1}{Vàng GLD: phần dư của mô hình tốt nhất.}{ch4-gold-resid}
\hinhpdf{ch4_gold_bars}{1}{Vàng GLD: RMSE và MAE theo kiến trúc và framework.}{ch4-gold-bars}

% =====================================================================
\clearpage
\section{So sánh tổng hợp và thí nghiệm bổ trợ}
\label{sec:rnn-ablation}

\subsection{Mô hình nào thắng dự báo ngây thơ?}

\hinhpdf{ch4_overview_ratio}{1}{Tỉ số RMSE mô hình / RMSE naive trên test cho 8 mô hình × 3 thực nghiệm.}{ch4-overview}

\nhanxet{Bức tranh rất rõ: trên e-commerce mọi RNN giảm sai số khoảng một phần ba so với naive (mô hình được
chọn theo validation đạt tỉ số \SL{ch4_summary/ecommerce/best_ratio} trên test), còn trên AAPL và vàng mọi mô hình nằm sát 1. Lý do không
nằm ở kiến trúc mà ở dữ liệu: doanh thu ngày có mức trung bình ổn định và nhiễu độc lập, nên mô hình nào học được mức
trung bình cũng thắng ``hôm qua''; giá tài sản là gần bước ngẫu nhiên, ``hôm qua'' đã là dự báo gần tối ưu.}

\subsection{PyTorch so với Keras}

\bang{ch4_fw_compare}{RMSE trên test của cùng kiến trúc trong hai framework}{ch4-fw}
\hinhpdf{ch4_cost}{1}{Số tham số (trái) và thời gian mỗi epoch (phải) của bốn kiến trúc trên AAPL.}{ch4-cost}

\nhanxet{Hai framework cho kết quả rất gần nhau; RMSE trung bình của bốn kiến trúc là
\SL{ch4_summary/ecommerce/torch_mean_rmse} (PyTorch) và \SL{ch4_summary/ecommerce/keras_mean_rmse} (Keras) trên
e-commerce, \SL{ch4_summary/aapl/torch_mean_rmse} và \SL{ch4_summary/aapl/keras_mean_rmse} trên AAPL,
\SL{ch4_summary/gold/torch_mean_rmse} và \SL{ch4_summary/gold/keras_mean_rmse} trên vàng. Chênh lệch nhỏ đến từ khởi tạo
mặc định khác nhau (Keras dùng Glorot và khởi tạo trực giao cho $W_{hh}$, đặt sẵn bias cổng quên LSTM bằng 1), cách
tính tham số (PyTorch có hai vector bias mỗi cổng) và tính ngẫu nhiên của huấn luyện. Ở quy mô dữ liệu này lựa chọn
framework không ảnh hưởng tới chất lượng; nó chỉ ảnh hưởng tới cách viết mã và tốc độ.}

\subsection{Chuẩn hoá Min--Max toàn cục so với chuẩn hoá tương đối}

Để kiểm chứng lập luận ở mục~\ref{sec:rnn-chuanhoa}, LSTM và GRU được huấn luyện lại với Min--Max toàn cục (fit trên
train, đúng cách làm của nhiều bài tham khảo) và so với chuẩn hoá tương đối.

\maNguon{ch4_data__Prepared._minmax}{Chuẩn hoá Min--Max toàn cục, chỉ dùng làm đối chứng}
\maNguon{ch4_extras__minmax_ablation}{Thí nghiệm đối chứng: cùng mô hình, hai cách chuẩn hoá}
\bang{ch4_minmax}{Ảnh hưởng của cách chuẩn hoá tới dự báo trên test}{ch4-minmax}
\hinhpdf{ch4_minmax}{1}{Dự báo của LSTM với Min--Max toàn cục (đỏ) và với chuẩn hoá tương đối (xanh, đứt) trên test
AAPL (trái) và vàng (phải).}{ch4-minmax}

\begin{phathien}{Phát hiện: $R^2$ âm của bài mẫu là lỗi chuẩn hoá, không phải lỗi mô hình}
Với Min--Max toàn cục, LSTM trên AAPL có RMSE \SL{ch4_extras/minmax/aapl/minmax_lstm/rmse} USD và
$R^2=$ \SL{ch4_extras/minmax/aapl/minmax_lstm/r2}; cùng mô hình với chuẩn hoá tương đối đạt
$R^2=$ \SL{ch4_extras/minmax/aapl/relative_lstm/r2}. Giá cao nhất của train là \SL{ch4_extras/minmax/aapl/train_max}
USD trong khi test lên tới \SL{ch4_extras/minmax/aapl/test_max}: đầu vào test nằm ngoài $[0,1]$, $\tanh$ bão hoà và dự
báo ``kẹt'' gần mức trần (Hình~\ref{fig:ch4-minmax}). Trên vàng, nơi giá test vẫn nằm phần lớn trong vùng train, khác
biệt nhỏ hơn: $R^2$ là \SL{ch4_extras/minmax/gold/minmax_lstm/r2} so với \SL{ch4_extras/minmax/gold/relative_lstm/r2}.
\end{phathien}

\subsection{Độ dài cửa sổ}

\maNguon{ch4_extras__window_ablation}{Thí nghiệm độ dài cửa sổ 10--90 bước}
\bang{ch4_window}{Tỉ số RMSE/naive theo độ dài cửa sổ}{ch4-window}
\hinhpdf{ch4_window}{1}{Tỉ số RMSE/naive của LSTM và GRU khi đổi độ dài cửa sổ.}{ch4-window-fig}

\nhanxet{Trên AAPL và vàng mọi độ dài cửa sổ cho tỉ số sát 1: thêm lịch sử không giúp vì lịch sử không chứa thông tin
về hướng. Trên e-commerce cửa sổ dài giúp ước lượng mức trung bình ổn định hơn; tốt nhất là LSTM với 60 ngày
(\SL{ch4_extras/window/ecommerce/lstm_w60/ratio}), nhưng cửa sổ dài cũng làm giảm số mẫu train của một chuỗi vốn đã
ngắn.}

\subsection{Dự báo nhiều bước bằng đệ quy}
\label{sec:rnn-horizon}

Ứng dụng web cần dự báo nhiều ngày tới. Ta dùng chiến lược đệ quy (Hình~\ref{fig:ch4-recursive}): dự báo một bước,
dựng hàng đặc trưng của ngày mới từ giá trị vừa dự báo, nối vào cửa sổ và lặp lại. Với giá, các cột giá cùng tài sản
(open, high, low, close) được dịch theo cùng tỉ lệ thay đổi, các cột khác giữ nguyên giá trị cuối; với e-commerce,
các đặc trưng phụ lấy trung bình 7 ngày gần nhất.

\maNguon{ch4_forecast__window_to_input}{Đưa một cửa sổ giá trị thô về đầu vào chuẩn hoá của mô hình}
\maNguon{ch4_forecast__next_row}{Dựng hàng đặc trưng của bước tiếp theo từ giá trị vừa dự báo}
\maNguon{ch4_forecast__recursive_forecast}{Dự báo đệ quy nhiều bước}
\maNguon{ch4_extras__horizon_eval}{Đánh giá sai số theo tầm dự báo trên nhiều gốc thời gian của tập test}

\bang{ch4_horizon}{RMSE theo tầm dự báo của LSTM đệ quy so với naive}{ch4-horizon}
\hinhpdf{ch4_horizon}{1}{RMSE theo tầm dự báo $h=1..10$ trên ba thực nghiệm.}{ch4-horizon-fig}

\nhanxet{Sai số tăng theo tầm dự báo ở cả ba bộ dữ liệu. Với AAPL, RMSE tăng từ \SL{ch4_extras/horizon/aapl/h1} ở $h=1$
lên \SL{ch4_extras/horizon/aapl/h10} ở $h=10$ (gấp \SL{ch4_extras/horizon/aapl/growth} lần), sát với naive
(\SL{ch4_extras/horizon/aapl/naive_h10}), đúng quy luật $\sqrt h$ của bước ngẫu nhiên. Trên e-commerce, sai số tăng
chậm hơn (\SL{ch4_extras/horizon/ecommerce/h1} lên \SL{ch4_extras/horizon/ecommerce/h10}) và vẫn thấp hơn naive
(\SL{ch4_extras/horizon/ecommerce/naive_h10}), vì dự báo đệ quy hội tụ về mức trung bình. Ứng dụng web vì vậy hiển thị
dự báo nhiều bước như một xu hướng tham khảo, không phải lời hứa về giá.}

% =====================================================================
\section{Triển khai hệ thống}
\label{sec:rnn-deploy}

Ba thực nghiệm được đóng gói thành mô-đun suy luận \texttt{tl3/serve\_ch4.py} mà Web App chung của tiểu luận gọi tới.
Mô-đun khai báo \texttt{SPECS} cho ba mã thực nghiệm \texttt{ch4\_ecommerce}, \texttt{ch4\_aapl}, \texttt{ch4\_gold};
mỗi mã có kiểu \texttt{series}, hai ô nhập (số bước dự báo 1--30 và chuỗi lịch sử tuỳ chọn), danh sách tám mô hình
\texttt{torch\_rnn} \dots\ \texttt{keras\_bilstm}, mặc định \texttt{torch\_lstm}. Hàm \texttt{predict} nạp mô hình lười
từ \texttt{models/ch4}, lấy 30 bước cuối của dữ liệu (hoặc chuỗi người dùng dán vào), dự báo bước kế tiếp và
\texttt{horizon} bước đệ quy, rồi trả về giá trị chính cùng lịch sử 90 bước, ngày tương ứng, chuỗi dự báo, ngày dự báo,
thay đổi phần trăm so với hôm nay, dự báo naive và chỉ số test của mô hình đã chọn để giao diện vẽ biểu đồ.

\maDai{serve_ch4__predict}{Hàm suy luận của Chương 4 cho Web App}

\begin{luuy}{Cách đọc dự báo trên giao diện}
Biểu đồ hiển thị lịch sử thực (màu đậm) và dự báo đệ quy (màu nhạt, nét đứt). Dựa trên kết quả mục~\ref{sec:rnn-ablation},
giao diện luôn hiện kèm dự báo naive và RMSE test để người dùng thấy mô hình có thực sự hơn ``giá hôm nay'' hay
không. Với AAPL và vàng, dự báo nhiều bước chỉ mang tính minh hoạ.
\end{luuy}

\subsection{Thực nghiệm 1 trên Web App}
\anhUI{ch4_ecommerce_web}{Web App: dự báo doanh thu ngày kế tiếp và 10 ngày tới của sàn thương mại điện tử.}{ch4-ecom-web}
\anhMobile{ch4_ecommerce_mobile}{Giao diện dự báo doanh thu trên điện thoại.}{ch4-ecom-mobile}

\subsection{Thực nghiệm 2 trên Web App}
\anhUI{ch4_aapl_web}{Web App: dự báo giá đóng cửa AAPL, so với naive và chỉ số test của mô hình.}{ch4-aapl-web}
\anhMobile{ch4_aapl_mobile}{Giao diện dự báo AAPL trên điện thoại.}{ch4-aapl-mobile}

\subsection{Thực nghiệm 3 trên Web App}
\anhUI{ch4_gold_web}{Web App: dự báo giá quỹ vàng GLD với lựa chọn tám mô hình.}{ch4-gold-web}
\anhMobile{ch4_gold_mobile}{Giao diện dự báo giá vàng trên điện thoại.}{ch4-gold-mobile}

% =====================================================================
\section{Kết luận chương}
\label{sec:rnn-ketluan}

Chương này đi từ lý thuyết đến triển khai mạng nơ-ron hồi quy:
\begin{itemize}
  \item Về lý thuyết, RNN là hệ động lực rời rạc có tham số dùng chung theo thời gian, huấn luyện bằng BPTT. Tích các
        Jacobian trong BPTT là nguồn gốc của gradient biến mất/bùng nổ, và đo trực tiếp cho thấy SimpleRNN mất gần
        hết gradient sau 50 bước, LSTM giữ được nhờ cổng quên gần 1, còn RNN-ReLU với $\rho>1$ bùng nổ.
  \item RNN cài đặt từ đầu bằng NumPy có gradient khớp sai phân số tới \SL{ch4_scratch/max_rel_err} và khớp
        \texttt{torch.nn.RNN}, rồi học được chuỗi sin có nhiễu tốt gần gấp đôi naive.
  \item Trên dữ liệu thực, chất lượng dự báo do \emph{dữ liệu} quyết định nhiều hơn \emph{kiến trúc}. Doanh thu ngày
        có ACF gần 0 nên mọi RNN thắng naive khoảng một phần ba; giá AAPL và vàng gần bước ngẫu nhiên nên mọi RNN chỉ
        ngang naive, độ chính xác hướng khoảng 50\%. SimpleRNN, LSTM, GRU, BiLSTM chênh nhau rất ít; PyTorch và Keras
        cũng vậy.
  \item Thiết kế tiền xử lý quan trọng hơn chọn mô hình: chuẩn hoá tương đối theo cửa sổ đưa $R^2$ của LSTM trên AAPL
        từ \SL{ch4_extras/minmax/aapl/minmax_lstm/r2} lên \SL{ch4_extras/minmax/aapl/relative_lstm/r2}. Luôn so với
        dự báo ngây thơ trước khi tin vào $R^2$.
  \item Ba mô hình được triển khai thành dịch vụ dự báo nhiều bước trong Web App, hiển thị kèm baseline để người dùng
        đọc kết quả đúng mức.
\end{itemize}
Hướng mở rộng tự nhiên gồm dự báo xác suất (khoảng tin cậy thay vì một điểm), mô hình hoá độ biến động thay vì hướng
giá, và so sánh với Transformer cho chuỗi thời gian hoặc các mô hình không gian trạng thái \cite{ch4:gu2023}.
